% Budget: 5 pages, ~1,500 words

\chapter{Methodology}\label{ch:methodology}

\section{Positionality}\label{sec:positionality}

As set out in the introduction (\ref{ch:introduction}), this project aims to map the different ways cults are conceived. I personally consider that I have entertained a certain fascination for this topic largely because of the many documentaries and media pieces that are available about these things, especially in the French context with l'Ordre du Temple Solaire, Raël and his followers, and a lot of other groups that have had entire documentaries made about them in the early 2000s. The main reason why I find this topic relevant currently is because I hold the intuition that the mechanisms at play in `cults' also play out in other social relationships - like the concept of `sectarian drifts' hint at - and that these mechanisms are also at play in the way people relate to new technologies, and more specifically to AI.

Because of this interest, the research actively seeks to gather as many different perspectives on the topic as possible, map them into a language space (the embedding space), and identify emergent patterns.

The underlying assumption is that many of the views would be contradictory and would enable us to better question the validity of the term `cult' as a meaningful category, as well as maybe find a more precise approach when trying to describe the phenomena emerging from these social relationships.

A further assumption at the start of this project is that, if we were to apply the criteria we use to assess whether a group belongs to the `cult' category, we would find that other groups we would not necessarily consider cults would also meet these criteria. This would suggest that the term `cult' is not a very precise category and is more of a social construct than a scientific one.

As discussed in the introduction (\ref{ch:introduction}), the resulting map shows how we conceive of cults through language, not the reality of cults.

I am aware that my own position, as a researcher and as a person, can influence the way people respond to my questions. The interviews were conducted with people from my own social and professional circles and with all people I had talked to at least once, some of whom were even aware of the objective of my ongoing research. There are many biases introduced by the fact that I was the one conducting the interviews and analysing them. The interviews were conducted in a semi-structured manner, with a set of questions I prepared in advance, while also allowing for flexibility and follow-up questions based on the interviewees' responses.

When conducting the interviews, I was aware that I was asking people to talk about a topic that can be sensitive for some. I aimed to be respectful and create a safe space for people to share their thoughts and experiences.

I also tried to be aware of my own biases and assumptions, avoid leading questions or imposing my views on the interviewees, and interrupt as little as possible.

Each of the three corpora brings its own biases, and the perspective of people who have been part of such groups is missing. These limitations are detailed in \ref{sec:limitations}.

\section{Data Collection}
% Here insert table of the 3 different epistemologies
\label{sec:DataCollection}
The goal is to find ways in which the concept of cult can be understood as a collective knowledge. This step constitutes the first step of the research project, and is the most important one. We're looking for associations, definitions, examples, oppositions, limits. In order to achieve this, three different types of epistemologies are gathered (Table~\ref{tab:epistemologies}).
The very first one that was considered was the one suggested by the MIVILUDES since it sparked this research in the first place and presents a ready-to-use ontology with its detailed modes of "sectarian drifts".
Since we focus on use in language, we are looking for the everyday understanding of the term. This was achieved through interviews with non-specialists. Finally, a more informed understanding of the term through researchers and scholars, which was achieved through a systematic literature review.

% Table requested by Celestin on 19 Nov (comment above); facts from notes/data_collection.md and notes/methodology.md
\begin{table}[htbp]
    \centering
    \scriptsize
    \renewcommand{\arraystretch}{1.3}
    \begin{tabularx}{\textwidth}{@{}>{\raggedright\arraybackslash}p{0.13\textwidth}*{3}{>{\raggedright\arraybackslash}X}@{}}
        \toprule
                      & MIVILUDES                                                                                                                       & Interviews                                                            & Literature                                                    \\
        \midrule
        Understanding & institutional, almost-legal                                                                                                     & everyday, non-specialist                                              & scholarly                                                     \\
        Source        & a French state body                                                                                                             & 28 people from my own circles                                         & researchers and scholars                                      \\
        Collection    & two documents: the 2022 to 2024 activity report and the page \foreignlanguage{french}{Comment identifier une dérive sectaire ?} & semi-structured interviews: 22 in person, 6 as Instagram text answers & systematic search in OpenAlex, screened and ranked            \\
        Material      & 17 criteria and 108 extracted expressions                                                                                       & 12,181 words; 430 participant segments after filler removal           & 68 documents (57 with full text); 5,741 extracted expressions \\
        Language      & French, with a working English translation                                                                                      & 22 in English, 6 in French                                            & English search terms                                          \\
        Period        & report 2022 to 2024; page accessed 27 Aug 2026                                                                                  & 5 Aug to 10 Sep 2026                                                  & published 2001 to 2024                                        \\
        \bottomrule
    \end{tabularx}
    \caption{The three corpora of the Data Collection.}
    \label{tab:epistemologies}
\end{table}

\subsection{MIVILUDES and Sectarian Drifts}

As mentioned in Section~\ref{sec:MIVILUDES}, a \textit{sectarian drift}\footnote{
    The English translation of the French noun 'Secte' IS 'Cult', so when the MIVILUDES refers to sectarian drifts, it is really talking about cult-like negative consequences. See Section~\ref{sec:defining_cult}.
}
can take different modalities, and do not all have to be fulfilled to be worthy of interest from the institutions. Therefore, within this specific framework, we should not conceive a cult as a group that fulfills a list of criteria (although this is the word they use), but rather as signs that they could be conceived as such.
Of course, this approach betrays the original intent of the MIVILUDES, which is indeed NOT to point at any object as a \textit{cult}, and only draw hints, a silhouette of what it could be. The MIVILUDES provide a rare attempt at an institutional, almost-legal, definition.
Another critical aspect of the sectarian-drift is that they don't look at the nature of the beliefs but only at behaviours or social relationships.
\subsubsection{The Documents}
Conveniently, the list of so-called \textit{criteria} is provided directly on their website\autocite[]{CommentIdentifierDerive} as 17 criteria and the list and translation can be found in Appendix~\ref{app:miviludes_criteria}.
Additionally, the MIVILUDES produces regular reports about its activities, which also constitute an interesting knowledge base, with statistics, analysis and writings by researchers.

\subsection{Interviews}
To get a different idea of the use of the word \textit{cult} in everyday language, it is also crucial to gather its use by non-specialists.

The interview methodology was designed to extract people's views without orienting the thought process of the participants, asking a very open-ended question (\enquote{When you hear the word cult, what comes to mind?}, see the protocol in Appendix~\ref{sec:interview_protocol}) and following a normal conversation flow, and only extract relevant data afterwards.
The open question has the advantage of letting the participant express themselves in their \textit{own words}, which is great for adding diversity of language, criteria, or examples.

% insert stats summary here
% Table requested by Celestin on 20 Nov (comment above); counts from thesis/corpus/interviews/metadata/database.csv
\begin{table}[htbp]
    \centering
    \scriptsize
    \renewcommand{\arraystretch}{1.3}
    \begin{tabularx}{0.85\textwidth}{@{}>{\raggedright\arraybackslash}p{0.18\textwidth}>{\raggedright\arraybackslash}X@{}}
        \toprule
        Participants & 28, all known to me: friends, friends of friends, acquaintances; Instagram stories as a second channel \\
        Gender       & 14 male, 12 female, 1 other, 1 redacted                                                                \\
        Age          & 26 to 97, median 34 (one redacted)                                                                     \\
        Residence    & Germany 18, France 6, Slovakia 1, Spain 1, United States 1, redacted 1                                 \\
        Method       & 22 in person, in four batches (11, 2, 7, 2); 6 Instagram text answers                                  \\
        Language     & 22 in English; 6 in French, in person, translated into English                                         \\
        Dates        & 5 Aug to 10 Sep 2026                                                                                   \\
        \bottomrule
    \end{tabularx}
    \caption{Interview participants.}
    \label{tab:participants}
\end{table}

The interviews were conducted spontaneously in different environments and transcribed in real time by Google Recorder\autocite{google_recorder_2026} for the audio. Then a manual cleanup was performed, first with an LLM\autocite{anthropic_2026} (see the notice of AI use, Section~\ref{sec:ai_notice}) and then manually against the audio. The LLM also performed a batching into one consistent format.

\begin{table}[htbp]
    \centering
    \scriptsize
    \renewcommand{\arraystretch}{1.3}
    \begin{tabularx}{0.85\textwidth}{@{}>{\raggedright\arraybackslash}p{0.22\textwidth}>{\raggedright\arraybackslash}X@{}}
        \toprule
        Words                   & 12,181 in total; median 287 per interview (111 to 1,536)                                                                                                               \\
        \quad in person         & 11,042; median 394 (156 to 1,536)                                                                                                                                      \\
        \quad Instagram         & 1,139; median 201 (111 to 273)                                                                                                                                         \\
        Turns                   & 145 questions, 188 answers                                                                                                                                             \\
        Participant segments    & 478; 430 after filler removal, median 13 per interview (2 to 44)                                                                                                       \\
        First example mentioned & 26 coded: 12 unclear, 5 classic new religious movement or religious group, 3 named group or movement, 2 metaphorical, 2 mixed, 1 digital or technology, 1 organisation \\
        Metadata kept           & age, gender, nationality, main language, residence, interview language, method, date                                                                                   \\
        \bottomrule
    \end{tabularx}
    \caption{Interview transcripts.}
    \label{tab:transcripts}
\end{table}

The final transcripts and statistics can be found in Appendix~\ref{app:interviews}.

\subsection{Literature}

OpenAlex
- An open, standard API, with citations and field labels

Search
- Fixed queries: cult, sect, new religious movements, extreme beliefs
- plus sectarian drifts, as the bridge to MIVILUDES
- Titles only, except for sectarian drifts (title and abstract)
- No NRM on its own: the acronym collides with other fields

Screening and ranking
- Screened by deduplication, type, field, subfield, topic, title patterns and dates
- Ranked by citations and by presence across queries, balanced by discipline
- Exclusion lists in the appendix

Result
- 68 documents: 33 articles, 23 books, 12 chapters, from 2001 to 2024
- 57 with full text
- 11 missing behind paywalls or out of print: an access bias
- Extracted with Docling, then checked by hand


\section{Embeddings}
